\documentclass[10pt,a4paper]{article}
\usepackage[margin=1.6cm]{geometry}
\usepackage[T1]{fontenc}
\usepackage{lmodern}
\usepackage{microtype}
\usepackage{booktabs}
\usepackage{multicol}
\usepackage{listings}
\usepackage{xcolor}
\usepackage[hidelinks]{hyperref}
\usepackage{titlesec}
\usepackage[inline]{enumitem}
\usepackage{float}

\setlist[itemize]{leftmargin=1.1em,itemsep=1pt,topsep=2pt,parsep=0pt}
\titlespacing{\section}{0pt}{5pt}{2pt}
\titleformat{\section}{\normalfont\bfseries\large}{\thesection.}{0.5em}{}
\titlespacing{\subsection}{0pt}{4pt}{1pt}
\titleformat{\subsection}{\normalfont\bfseries}{}{0em}{}
\setlength{\parindent}{0pt}
\setlength{\columnsep}{0.8cm}
\pagestyle{empty}

\lstdefinestyle{sparql}{
  basicstyle=\ttfamily\scriptsize,
  keywordstyle=\color{blue!60!black}\bfseries,
  commentstyle=\color{black!55}\itshape,
  morekeywords={PREFIX,SELECT,WHERE,VALUES,OPTIONAL,ORDER,BY,FILTER,a},
  morecomment=[l]{\#},
  columns=fullflexible,
  xleftmargin=2pt,
  aboveskip=3pt,belowskip=3pt,
  breaklines=true,
}

\title{\vspace{-2cm}\textbf{Nathan's Assignment Documentation}}
% \author{Nathan Schneider Gavenski}
% \author{\vspace{-0.4cm}}
\date{\vspace{-1cm}}

\begin{document}
\maketitle
\vspace{-0.8cm}

\begin{multicols}{2}

\section{Scope}

AIAct-Onto is an OWL~2 DL ontology of Regulation (EU) 2024/1689.
It is built as a vertical slice instead of a complete formalisation.
The selected scope is as follows:
\begin{enumerate*}[label=\textbf{(\arabic*)},itemjoin={{; }},itemjoin*={{; and }}]
      \item Art.~3 (definitions)
      \item Art.~5 (prohibited practices)
      \item Art.~6 and Annexes~I and~III (high-risk classification)
      \item Art.~8--17 (requirements and provider obligations)
      \item Art.~26--27 (deployer obligations)
      \item Art.~50 (transparency)
      \item Art.~51--55 (general-purpose AI models)
      \item Art.~64 and~70 (authorities).
\end{enumerate*}
% In scope: Art.~3 (definitions), Art.~5 (prohibited practices), Art.~6 and Annexes~I and~III (high-risk classification), Art.~8--17 (requirements and provider obligations), Art.~26--27 (deployer obligations), Art.~50 (transparency), Art.~51--55 (general-purpose AI models), and Art.~64 and~70 (authorities): 
Yielding a total of $23$ articles and $2$ annexes.

The released ontology holds $69$ classes, $160$ individuals, $11$ object properties and $1{,}709$ axioms ($471$ logical).
$134$ of its entities came from the pipeline; the rest is the hand-authored skeleton.
Coverage by type: $62$ obligations, $32$ AI system categories, $13$ requirements, $8$ prohibited practices, $7$ high-risk areas, $7$ actor roles and $5$ authorities.

The ontology reuses five existing vocabularies.
\emph{AI Risk Ontology} (AIRO) models AI systems and the roles around them, so \texttt{aiact:AISystem}, \texttt{ProviderRole} and \texttt{DeployerRole} map to \texttt{airo:AISystem}, \texttt{airo:AIProvider} and \texttt{airo:AIDeployer}.
\emph{Vocabulary of AI Risks} (VAIR), a companion to AIRO, supplies taxonomies of AI capabilities and application domains, which the Annex~III areas map onto.
\emph{Data Privacy Vocabulary} (DPV), a W3C standard, models organisations, roles and duties in a GDPR setting, close to how the Act treats operators, so \texttt{Organisation} and \texttt{Obligation} map to their DPV counterparts.
\emph{European Legislation Identifier} (ELI) is the EU's own scheme for citing legislation: every provision here is identified by its ELI URI ($21$ linked) and \texttt{aiact:LegalProvision} is a subclass of \texttt{eli:LegalResourceSubdivision}, so each annotation resolves on EUR-Lex, and is not a bare string.
\emph{Simple Knowledge Organization System} (SKOS), a W3C standard, carries the one-sentence definition of every entity through \texttt{skos:definition}.
\emph{Provenance Ontology} (PROV-O), a W3C standard, records which pipeline run generated each entity, through \texttt{prov:wasGeneratedBy} pointing at a \texttt{prov:Activity}, so every axiom traces back to its run, model and curation decision.

Each design choice below is recorded in the repository \texttt{README} as a numbered architecture decision record, stating the choice and why it was made, and govern the mapping style.
The mappings to AIRO, VAIR and DPV use \texttt{skos:closeMatch}, not \texttt{rdfs:subClassOf} or \texttt{owl:equivalentClass}.
A SKOS mapping records the correspondence without importing the other ontology's axioms, so nothing there can make this one inconsistent.
Subsumption is used only where the semantics are certain, as with ELI.
All mappings live in a separate file, so they can be included or left out.

\section{The AI-assisted pipeline}

Nine file-based stages, each rerunnable in isolation, driven by \texttt{make}. 
Both model stages run \emph{Qwen 3.5 (4B)} at temperature $0$.
This avoids API costs and rate limits; any other LLM would do, with some prompt tuning.

\textbf{S1 Ingest} snapshots the EUR-Lex HTML and records its SHA-256.
\textbf{S2 Chunk} splits it into one JSON file per article or annex, keeping paragraph and point
numbering so references reach point level (\texttt{5(1)(c)}).
\textbf{S3 Extract} sends one chunk per call, with every passage prefixed by its reference, and
constrains the reply to a schema (term, type, label, definition, article reference,
verbatim quote, bearer, applies-to, confidence). 181 candidate records.
\textbf{S4 Verify} applies a deterministic grounding check, then an LLM critic, then deduplication:
160 records.
\textbf{S5 Curate} is a \texttt{typer}/\texttt{rich} review tool: accept, edit, reject or merge, with every decision appended to a log.
It allows for a human-in-the-loop approach, yielding a total of 134 entities.
\textbf{S6 Build} turns the curated records into OWL.
\textbf{S7 Validate} runs the reasoner, the DL profile check, SHACL and ROBOT's quality report.
\textbf{S8--S9} produce metrics and documentation.
Finally, we take some precautions to avoid accuracy pitfalls:

\begin{itemize}
  \item \textbf{No model writes OWL.} It only fills a pre-defined JSON schema; all axioms are
        produced by deterministic code, thus the build is reproducible from the curated data.
  \item \textbf{Grounding by construction.} Every record carries a verbatim quote, and a non-LLM
        string check rejects any quote that is not in the source (fuzzy ratio $\geq 95$ after
        normalising whitespace and quote characters). 179 of 181 passed; 158 of the survivors matched
        the source exactly.
  \item \textbf{Deterministic checks before model judgement}: the critic only judges records
        that are already grounded.
  \item \textbf{Provenance throughout the pipeline.} Each entity records its article reference, its quote, the ELI
        URI of its provision, its curation status and the run that generated it.
\end{itemize}

\section{Key modelling decisions}

\begin{itemize}
  \item \textbf{Roles instead of agent subclasses}. One organisation can be a provider for one
        system and a deployer for another, so \texttt{ProviderRole} is something an
        \texttt{Organisation} holds, as in AIRO and DPV.
  \item \textbf{Obligations and prohibited practices as individuals}. They are enumerated
        items in specific articles and not open categories. Classes are used only where the reasoner
        must classify, namely the risk tiers.
  \item \textbf{Requirement $\neq$ Obligation.} Art.~8--15 state properties a high-risk system must
        have; Art.~16(a) separately obliges the provider to ensure compliance with them. Collapsing
        the two would make Art.~16(a) circular.
  \item \textbf{Risk tiers are defined.} \texttt{ProhibitedAISystem} $\equiv$
        \texttt{AISystem} $\sqcap$ $\exists$\texttt{uses\-Practice}.\allowbreak\texttt{Prohibited\-Practice},
        and likewise \texttt{AnnexIII\-HighRisk\-AISystem} over \texttt{usedInArea}, so the reasoner does the
        classification. The four tiers are pairwise disjoint. The Art.~6(3) derogation is
        deliberately ignored.
  \item \textbf{Punning} lets an obligation point at a class of system or a role
        (\texttt{appliesToSystem}, \texttt{obligationOf}) while staying in OWL~2 DL.
\end{itemize}

\section{Competency questions}

We provide six competency questions (CQ) that the ontology can answer, which are the ones the Act itself is designed to answer:
\begin{enumerate*}[label=\textbf{CQ\arabic*},itemjoin={{; }},itemjoin*={{; and }}]
  \item Which obligations does a provider of a high-risk AI system hold, and from which article?
  \item Which practices does the Act prohibit, and in which words?
  \item Which areas make a system high-risk?
  \item What must a high-risk AI system satisfy?
  \item For general-purpose models, which duties fall on the provider and which on the authorised
        representative?
  \item Which systems are prohibited or high-risk, and why?
\end{enumerate*}

All six SPARQL queries are executable (and located in \texttt{queries/*.rq}). 
As an example, CQ1 is shown below. It returns every obligation a provider holds for a high-risk AI
system, with its label, article reference and definition, ordered by article.
\lstinputlisting[style=sparql,firstline=6]{../queries/cq1_provider_obligations.rq}

CQ6 is the only one that needs the reasoner.
No example is asserted to be high-risk or prohibited, so CQ6 returns nothing without a reasoner and two rows with one, which is why it makes a test of the axioms.
% \lstinputlisting[style=sparql,firstline=6]{../queries/cq6_classified_systems.rq}

\section{Validation}

The classification test is the substantive one: a fictional CV screening tool asserted only as an AI
system used in the employment area is inferred to be an Annex~III high-risk system; a social scoring
platform is inferred to be prohibited; a chat assistant stays unclassified. Feeding the reasoner a
system that is both prohibited and high-risk makes it report inconsistency, so the disjointness
axioms are doing work.
\vspace{-0.5cm}
\begin{table}[H]
      \centering
      \scriptsize
      \begin{tabular*}{0.9\columnwidth}{@{\quad}l@{\extracolsep\fill}l@{\qquad}}
            \toprule
            \textbf{Check} & \textbf{Result} \\
            \midrule
            Parses as OWL & 1{,}774 triples \\
            OWL 2 DL profile & in profile \\
            Consistency (HermiT) & consistent, 0 unsatisfiable \\
            Classification (HermiT) & 3 of 3 as expected \\
            Annotation completeness & SHACL conforms \\
            ROBOT report & 0 errors \\
            \bottomrule
      \end{tabular*}
\end{table}
\vspace{-0.6cm}
Getting these checks to pass found five real defects, including two undeclared annotation properties
that put the ontology outside OWL~2 DL, and six generated IRIs that collided with hand-authored
classes and produced \texttt{AISystem} $\sqsubseteq$ \texttt{AISystem}.

\section{What the review step changed}

Since the curation step would be quite lengthy to execute for this assignment, we use an AI assistant to simulate the human review step.
In a paper or real deployment scenario, a human would be the reviewer, guaranteeing that the curated ontology is faithful and no hallucinations are introduced.
Of the 160 verified records, 76 were accepted unchanged, 58 were edited, 20 were merged into another entity and 6 were rejected outright.
\vskip -0.5em
\begin{table}[H]
      \centering
      \scriptsize
      \begin{tabular*}{0.9\columnwidth}{@{\quad}l@{\extracolsep\fill}rrrr@{\quad}}
            \toprule
            \textbf{Critic verdict} & \textbf{acc.} & \textbf{edit} & \textbf{merge} & \textbf{rej.} \\
            \midrule
            supported (44)   & 31 & 10 & 2  & 1 \\
            partially (113)  & 44 & 47 & 18 & 4 \\
            unsupported (3)  & 1  & 1  & 0  & 1 \\
            \bottomrule
      \end{tabular*}
\end{table}
\vspace{-0.4cm}
The critic's verdict predicted the outcome poorly: 30\% of what it called supported was still
changed, and 39\% of what it called partially supported was accepted unchanged.
Four findings stand out.

\textbf{(a) The critic pushed a systematic misclassification.} In 35 records it read ``shall'' in
Art.~9--15 and concluded the entity type was wrong and these were provider obligations. Following it
would have erased the Requirement/Obligation distinction the ontology is built on, and left Art.~16(a)
as a circular duty: an obligation to comply with obligations.

\textbf{(b) The critic invented a fact.} For Art.~51(2) it claimed the compute threshold was
$10^{28}$ and the extracted $10^{25}$ was wrong. The Act says $10^{25}$. A reviewer trusting the
critic would have introduced an error no downstream check could catch, because the value sits in a
paraphrased definition, not a quoted span. Two further verdicts claiming that ``the record added
this phrase'' were fabrications of the same kind.

\textbf{(c) No automated check scores completeness.} Grounding and the critic both ask whether a
record is faithful. 29 obligations arrived with no bearer and 40 records with no applies-to value, so \texttt{obligationOf} could not be asserted and CQ1 would have returned nothing. After review, all 62 obligations carry a bearer.

\textbf{(d) The critic cannot see across records.} It marked ``Special Categories Processing
Conditions'' \emph{supported}: a faithful definition with a verbatim quote. It had no way to see that
S4's deduplication had collapsed the six distinct conditions of Art.~10(5) into one entity.

\section{Metrics and proposed evaluation}

Grounding pass rate $98.9\%$, of which $98.8\%$ matched the source exactly.
Annotation completeness $100\%$: all $157$ entities carry a label and a definition, and every curated
entity also carries an article reference, a quote and an ELI link.
Bearer coverage $62/62$ obligations.
Hierarchy depth $2$, with a mean of $7.1$ children per parent.
$8$ of $30$ core classes are aligned to an external ontology.
The ontology is consistent with $0$ unsatisfiable classes, and answers $6/6$ competency questions.
Critic and reviewer agree on $60.6\%$ of records, Cohen's $\kappa = 0.25$, a fair agreement that
quantifies how little the critic's verdict predicts the review outcome.

A fuller evaluation would need: \textbf{(i)} a gold standard, 3--5 articles annotated by hand before
running the pipeline, giving precision and recall per entity type: the most defensible single
accuracy number, and the main gap here; \textbf{(ii)} inter-annotator agreement between two
independent reviewers, separating the critic's weakness from the difficulty of the judgement;
\textbf{(iii)} curation effort per record, which decides whether the approach scales;
\textbf{(iv)} OOPS! for modelling pitfalls, and a task-based evaluation answering a compliance
question with and without the ontology.

\textbf{Known limitations.} The reviewer was an AI, so it shares the extractor's blind spots. The
alignment IRIs are unverified against downloaded copies of AIRO, VAIR and DPV. The Art.~6(3)
derogation, penalties and conformity assessment bodies are out of scope. A 4B model on a laptop GPU
sets the quality ceiling; a larger one needs no code change.

\end{multicols}
\end{document}
